% Chapter 7 -- Maintaining Turbo Initech.  Sources (Law 1): os/tps/tps.pas
% (structure: constants and fixed arenas; the lexer; the parse routines,
% one per production, shared by the syntax, semantic and generation
% passes through ParseFailed/TypeMode/GenFlags; the driver's main block);
% os/tps/CORPUS.md (fixture + hand golden, MODE DELPHI/B+/H- pins, forward
% discipline), CODEGEN.md s9 (the 72 KiB BSS reserve and the 6F000h image
% ceiling), the Makefile corpus list TPS_COMPILER_CORPUS; ADR-0007 C-2
% (deferred families need an ADR amendment); pubs/turbo.mk test-pubs-turbo (the
% manuals are checked against the compiler).
\Gates{test-compiler, test-compiler-os, test-tps-lex-build, test-pubs-turbo}
\chapter{Maintaining Turbo Initech}
\ix{maintaining Turbo Initech}

This chapter is for the programmer who changes Turbo Initech: what the
source looks like, which rules a change must keep, and what has to change
with it.

\section{The shape of the source}
\ix{tps.pas}
The whole compiler is the one file \file{os/tps/tps.pas}, in this order:

\begin{enumerate}
\item the constants: the capacities, the type, symbol and token codes, the
  error codes and the names of the grammar rules;
\item the global variables: the source window, the name pool, the symbol
  table, the scratch tables for headings and calls, and the state of the
  code generator, every one a fixed array or a simple variable;
\item the lexer: reading the source a block at a time, skipping blanks and
  comments, and recognizing each kind of token;
\item the parser: one routine for each rule of the grammar, named after it
  (\id{ParseFactor}, \id{ParseStatement}, \dots);
\item the four passes, \id{RunLexer}, \id{RunParser}, \id{RunTypeChecker} and
  \id{RunCodeGenerator};
\item the main block, which runs the passes in order and prints the lines
  between them.
\end{enumerate}

The same parsing routines serve the syntax, semantic and generation
passes: a pass that checks types sets \id{TypeMode}, and the generation pass
sets the first of \id{GenFlags}; each routine checks or writes code as the
mode asks while it reads.  A change to the grammar is therefore one change,
seen by all three passes.

\section{Rules a change must keep}
\begin{itemize}
\item \textbf{Both compilers accept the source.}  The genesis compiler must
  still compile \file{tps.pas}: declare every routine before its first call,
  using \rw{forward}, and do not use \rw{case} as a name.
\item \textbf{The language stays as defined.}  Pointers, sets, enumerated
  types, units, real numbers, objects and inline assembly are left out on
  purpose; adding any of them is a change to the architecture of Turbo
  Initech, decided before it is made, not in a change to the source.
\item \textbf{Guards are nested.}  Both operands of \rw{and} and \rw{or} are
  evaluated, in the compiler's own source as in any program.
\item \textbf{Tables stay in order.}  New tables are arrays kept in the
  order of the declarations; no table may be ordered by anything else.
\item \textbf{Capacities fail loudly.}  Every table has a fixed size, and
  running out of room must report an error that names the limit, never
  truncate or wrap around.
\item \textbf{The compiler fits.}  The build checks that the compiler's
  global variables stay within their reserve and that its image ends below
  6F000h (procedure TIN-P-02).
\end{itemize}

\section{What changes with a change}
\SideHead{Tests}A new feature of the language comes with a corpus program
whose expected output is worked out by hand and stored beside it; the
program is added to the corpus list of the build and to Appendix A, and the
Free Pascal compiler must agree with the hand-worked output.

\SideHead{Manuals}The Turbo Initech manuals are checked against the
compiler every time they are built.  A new reserved word, routine, grammar
rule, message, output line or capacity that the manuals do not describe,
or a description of something the compiler no longer has, stops the check;
so does any example whose output changes.  A change to the compiler
therefore carries the change to the manuals.

\SideHead{Acceptance}A changed compiler is accepted by the whole procedure
of Chapter 4.
